\section{A sharp refinement certificate}\label{sec:certificate}
Let $\widehat K_j(x,\cdot)\in\Ptwo(E_j)$ be learned detail kernels and let $\widehat\mu_j$ be their sequentially generated laws. The inputs to our bound are a conditional fitting error and a learned context sensitivity:
\begin{align}\label{eq:kernel-errors}
 \eps_j^2&=\int_{H_{j-1}}W_2^2(K_j(x,\cdot),\widehat K_j(x,\cdot))\,\mu_{j-1}(\mathrm dx),\\
 W_2(\widehat K_j(x,\cdot),\widehat K_j(y,\cdot))&\le L_j\norm{x-y}.
 \label{eq:context-lipschitz}
\end{align}
For $j=0$, set $\mu_{-1}=\delta_0$, $L_0=0$. Require \eqref{eq:context-lipschitz} on all contexts needed by a coupling of the true and learned recursions; a global bound on $H_{j-1}$ is a sufficient convenient assumption. The errors in \eqref{eq:kernel-errors} use the \emph{target} context distribution. The sensitivity bound controls how the learned kernel changes when it is fed generated contexts instead.

\begin{theorem}[Orthogonal refinement certificate]\label{thm:certificate}
Assume \eqref{eq:kernel-errors} is finite and \eqref{eq:context-lipschitz} holds through level $m$. Define
\begin{equation}\label{eq:sharp-recurrence}
 d_{-1}=0,\qquad d_j^2=d_{j-1}^2+(\eps_j+L_jd_{j-1})^2.
\end{equation}
Then
\begin{equation}\label{eq:finite-certificate}
 W_2(\mu_j,\widehat\mu_j)\le d_j\qquad (0\le j\le m).
\end{equation}
Writing $\tau_m^2=\E\norm{(I-P_m)U}^2$, the exact identity and its bound are
\begin{equation}\label{eq:pythagorean-certificate}
 W_2^2(\mu,\widehat\mu_m)
 =\tau_m^2+W_2^2(\mu_m,\widehat\mu_m)
 \le\tau_m^2+d_m^2.
\end{equation}
In particular,
\begin{align}\label{eq:sum-certificate}
 d_m^2&\le 2\sum_{j=0}^m\eps_j^2
                \prod_{k=j+1}^m(1+2L_k^2)\notag\\
 &\le 2\exp\left(2\sum_{k=1}^mL_k^2\right)\sum_{j=0}^m\eps_j^2.
\end{align}
If $\sum_jL_j^2<\infty$ and $\sum_j\eps_j^2<\infty$, the learned tower converges almost surely and in mean square to an $H$-valued random element with law $\widehat\mu$, and
\begin{equation}\label{eq:infinite-certificate}
 W_2(\mu,\widehat\mu)
 \le \lim_{m\to\infty}d_m
 \le\sqrt2\exp\left(\sum_{j\ge1}L_j^2\right)
             \left(\sum_{j\ge0}\eps_j^2\right)^{1/2}.
\end{equation}
\end{theorem}
\begin{proof}
Suppose $(X,\widehat X)$ couples $\mu_{j-1}$ and $\widehat\mu_{j-1}$ with $L^2$ distance $D$. Conditional on $(X,\widehat X)=(x,y)$, couple three detail variables with marginal laws
\[
 B\sim K_j(x,\cdot),\quad
 \widetilde B\sim\widehat K_j(x,\cdot),\quad
 \widehat B\sim\widehat K_j(y,\cdot).
\]
Choose optimal couplings for the first two and last two laws and glue them along $\widetilde B$. For the continuous nonnegative cost $\|b-b'\|^2$ on the Polish space $E_j\times E_j$, \citet[Corollary~5.22]{villani2009} supplies Borel selections of optimal plans as measurable functions of their marginals. The target and learned kernels are Borel maps into $\mathcal P_2(E_j)$, so these selections remain measurable in $(x,y)$. Parameterized disintegration of the two plans along their common middle marginal gives a measurable gluing kernel. Its successive sampling defines a genuine joint probability law; the same selection argument can include an external conditioning parameter. Minkowski's inequality and \eqref{eq:kernel-errors}--\eqref{eq:context-lipschitz} give
\begin{equation}\label{eq:detail-coupling-bound}
 (\E\norm{B-\widehat B}^2)^{1/2}
 \le\eps_j+L_jD.
\end{equation}
The full discrepancy lies in orthogonal old and new bands, so
\[
 \E\norm{(X+B)-(\widehat X+\widehat B)}^2
 =D^2+\E\norm{B-\widehat B}^2
 \le D^2+(\eps_j+L_jD)^2.
\]
Starting with equal zero contexts and using this construction recursively proves \eqref{eq:finite-certificate}. Finite second moments of the learned outputs also follow inductively: a target-averaged finite fitting error gives finite moments at target contexts, and the Lipschitz bound transfers them to the generated contexts.

For \eqref{eq:pythagorean-certificate}, any coupling $(U,Y)$ with $Y\in H_m$ satisfies
\[
 \norm{U-Y}^2=\norm{(I-P_m)U}^2+\norm{P_mU-Y}^2.
\]
Taking infima gives the lower inequality in the claimed identity. Conversely, lift an optimal coupling of $(P_m)\push\mu$ and $\widehat\mu_m$ by sampling $U$ from its conditional law given $P_mU$. This realizes the sum of the two optimal terms and proves equality.

Using $(a+b)^2\le2a^2+2b^2$ in \eqref{eq:sharp-recurrence} gives
\[
 d_j^2\le(1+2L_j^2)d_{j-1}^2+2\eps_j^2.
\]
Iteration and $1+x\le e^x$ prove \eqref{eq:sum-certificate}. Under the infinite summability assumptions, $d_m$ is bounded. The triangle inequality for root second moments yields
\[
 (\E\norm{\widehat X_m}^2)^{1/2}
 \le (\E\norm{P_mU}^2)^{1/2}+W_2(\mu_m,\widehat\mu_m)
 \le(\E\norm U^2)^{1/2}+\sup_m d_m.
\]
The learned energy condition in \Cref{thm:representation} follows. Hence $\widehat\mu_m\to\widehat\mu$ in $W_2$, while $\mu_m\to\mu$ in $W_2$. Taking limits in \eqref{eq:finite-certificate} proves \eqref{eq:infinite-certificate}.
\end{proof}

\subsection{What orthogonality adds to kernel perturbation theory}\label{sec:comparison}
The usual triangle-inequality propagation for the appended-state kernel has the comparison update
\[
 b_j=\sqrt{1+L_j^2}\,b_{j-1}+\eps_j.
\]
Indeed, the learned transition $x\mapsto(x,\widehat B_j(x))$ has Wasserstein sensitivity at most $\sqrt{1+L_j^2}$. Treating fitting as a subsequent unrestricted perturbation then gives the displayed bound, in the spirit of Wasserstein kernel perturbation arguments \citep{rudolf2018}. Our proof instead retains the direction of the new error:
\begin{equation}\label{eq:orthogonal-improvement}
 D^2+(\eps+LD)^2\le
 \bigl(\sqrt{1+L^2}D+\eps\bigr)^2,
 \qquad D,L,\eps\ge0.
\end{equation}
The difference of the right and left sides is
$2\eps D(\sqrt{1+L^2}-L)\ge0$.
When $L_j=0$, the generic update accumulates $\sum_j\eps_j$, while the exact orthogonal update accumulates $(\sum_j\eps_j^2)^{1/2}$. With $\eps_j=1/(j+1)$, only the latter remains finite. This is the precise gain: persistence of old coordinates and orthogonality of newly fitted detail improve how perturbations accumulate. Disintegration, optimal-coupling selection, and the underlying Wasserstein inequalities are classical ingredients.

\begin{proposition}[Sharpness of the recurrence]\label{prop:sharpness}
For any nonnegative finite sequences $(\eps_j)$ and $(L_j)$ with $L_0=0$, equality in \eqref{eq:finite-certificate} can hold at every level.
\end{proposition}
\begin{proof}
Take scalar orthogonal bands, target $\mu=\delta_0$, and $K_j(x,\cdot)=\delta_0$. Define
\[
 \widehat K_j(x,\cdot)=\delta_{\eps_j+L_j\norm{x}}.
\]
This kernel is $L_j$-Lipschitz in $W_2$, and its error at the only target context $x=0$ is $\eps_j$. The generated vector is deterministic. If its norm before refinement is $d_{j-1}$, the new coordinate is $\eps_j+L_jd_{j-1}$, giving \eqref{eq:sharp-recurrence} exactly. Its distance from $\delta_0$ is its norm.
\end{proof}

The sharpness statement concerns the general kernel theorem, which includes deterministic kernels. It does not assert equality within the regular Gaussian-reference ODE subclass. The exponential simplification \eqref{eq:sum-certificate} is convenient but not sharp; for context-independent kernels $L_j=0$, the exact recurrence gives $d_m^2=\sum_{j=0}^m\eps_j^2$.

\subsection{The exact worst-case energy gain}\label{sec:optimal-gain}
The individual recurrence is sharp, but its exponential simplification can be loose. We can characterize the best constant when only the aggregate conditional fitting energy is available. This also identifies precisely when a resolution-uniform guarantee is possible over the full kernel class.

For fixed nonnegative sensitivities $(L_j)_{j\ge1}$, define $g_0=1$ and
\begin{align}\label{eq:optimal-gain}
 A_j&=(1+L_j^2)g_{j-1},\\
 g_j&=\frac{A_j+1+\sqrt{(A_j-1)^2+4L_j^2g_{j-1}}}{2}.
 \notag
\end{align}
The index $0$ here denotes the first band, as in \Cref{thm:certificate}.

\begin{theorem}[Optimal energy amplification and its threshold]\label{thm:optimal-gain}
For every nonnegative error vector $(\eps_0,\ldots,\eps_m)$, the exact recurrence \eqref{eq:sharp-recurrence} satisfies
\begin{equation}\label{eq:optimal-energy-bound}
 d_m^2\le g_m\sum_{j=0}^m\eps_j^2.
\end{equation}
The factor $g_m$ is the smallest possible uniform constant. Equality is attained by a nonnegative error vector of unit squared sum and by the deterministic kernels in \Cref{prop:sharpness}. Consequently it is also the optimal worst-case constant for the general conditional-kernel theorem, given only those sensitivity bounds and aggregate fitting energy.

Moreover,
\begin{equation}\label{eq:gain-threshold}
 \sup_m g_m<\infty
 \quad\Longleftrightarrow\quad
 \sum_{j\ge1}L_j^2<\infty.
\end{equation}
Thus squared-sensitivity summability is necessary and sufficient for a uniform robustness guarantee over this class. It remains a sufficient, not necessary, condition for the accuracy or existence of any one particular model.
\end{theorem}
\begin{proof}
At the first band $d_0^2=\eps_0^2$, so $g_0=1$ is exact. Suppose the assertion holds through $j-1$. Write
$u=(\sum_{k<j}\eps_k^2)^{1/2}$ and $v=\eps_j$. Monotonicity of the recurrence bounds the next squared radius by
\[
 g_{j-1}u^2+(v+L_j\sqrt{g_{j-1}}u)^2
 =\begin{pmatrix}u&v\end{pmatrix}
 \begin{pmatrix}
 (1+L_j^2)g_{j-1}&L_j\sqrt{g_{j-1}}\\
 L_j\sqrt{g_{j-1}}&1
 \end{pmatrix}
 \begin{pmatrix}u\\v\end{pmatrix}.
\]
Its largest eigenvalue is exactly \eqref{eq:optimal-gain}. The spectral bound gives \eqref{eq:optimal-energy-bound}. This argument uses a two-dimensional matrix only; the coordinate dimension of the generated field is irrelevant.

For sharpness choose a unit eigenvector for the largest eigenvalue with nonnegative coordinates $(u,v)$. Such a choice exists because the off-diagonal entry is nonnegative: replacing either coordinate by its absolute value does not lower the quadratic form. Scale an attaining error vector from the previous step by $u$ and append $v$. The recurrence is homogeneous, so the resulting vector has unit squared sum and attains $g_j$. When the off-diagonal entry is zero, the larger diagonal coordinate gives the same conclusion, including the degenerate identity matrix. Induction and \Cref{prop:sharpness} establish attainment by actual kernels.

If the sensitivity square sum is finite, \eqref{eq:sum-certificate} gives
$g_m\le2\exp(2\sum_{j\ge1}L_j^2)$, uniformly in $m$. Conversely, take $\eps_0=1$ and all later fitting errors zero. The exact recurrence gives
\begin{equation}\label{eq:seed-product}
 d_m^2=\prod_{j=1}^m(1+L_j^2)\le g_m.
\end{equation}
This product diverges when $\sum_jL_j^2=\infty$: either infinitely many squared sensitivities exceed one, or eventually
$\log(1+L_j^2)\ge L_j^2/2$.
This proves necessity. In the attaining deterministic construction, the diverging product also prevents a Hilbert-valued finite-energy limiting output, even though the entire fitting error is confined to the first band.
\end{proof}

Equation \eqref{eq:optimal-gain} is useful when validation supplies only a total fitting-energy budget. With separate bandwise errors, the original recurrence can be smaller and should be retained. The gain theorem concerns unrestricted conditional kernels; equality need not be achievable by the smooth neural family of \Cref{sec:spectral-network}. It gives an exact stability threshold for the abstract class rather than a new convergence claim for every trained architecture.

The program \texttt{experiments/optimal\_gain.py} reconstructs the attaining error vector from the successive two-dimensional eigenvectors and verifies its recurrence directly. It also compares $\sqrt{g_m}$ with the exponential majorant for three sensitivity schedules; the numerical table reports finite levels, while \eqref{eq:gain-threshold} supplies the infinite conclusion.
\begin{table}[htbp]
\centering\small
\begin{tabular}{rrrr}
\toprule
$p$ & Last band $m$ & Optimal gain & Exponential gain \\
\midrule
1.00 & 8 & 1.4596 & 2.0718 \\
1.00 & 32 & 1.4891 & 2.1172 \\
1.00 & 128 & 1.4969 & 2.1294 \\
0.50 & 8 & 1.8029 & 2.7900 \\
0.50 & 32 & 2.2464 & 3.9009 \\
0.50 & 128 & 2.7563 & 5.5006 \\
0.25 & 8 & 2.2705 & 4.2183 \\
0.25 & 32 & 4.7895 & 16.9786 \\
0.25 & 128 & 19.5679 & 284.1154 \\
\bottomrule
\end{tabular}

\caption{Exact aggregate-error gains for $L_j=0.5j^{-p}$. Each entry compares the optimal gain $\sqrt{g_m}$ with the sufficient exponential gain $\sqrt2\exp(\sum_{j=1}^mL_j^2)$. Only $p>1/2$ admits a bounded infinite-hierarchy gain.}\label{tab:optimal-gain}
\end{table}

\FloatBarrier

\begin{corollary}[Convergence of improving model families]\label{cor:learning-consistency}
Let a sequence of coherent learned towers be indexed by $n$, with
\[
 \sup_n\sum_jL_{n,j}^2<\infty,\qquad \sum_j\eps_{n,j}^2\longrightarrow0.
\]
Their infinite laws converge to $\mu$ in $W_2$. If only level $m_n$ is sampled and $\tau_{m_n}\to0$, the finite generated laws also converge to $\mu$ in $W_2$.
\end{corollary}
\begin{proof}
Apply \eqref{eq:infinite-certificate} and \eqref{eq:pythagorean-certificate} with the uniform sensitivity bound.
\end{proof}

No finite-data convergence rate is asserted here. The corollary identifies the aggregate conditional error and stability conditions that a statistical learning theorem would need to supply. For a specified neural class, \Cref{cor:validated-law} supplies a separate finite-sample audit of these error inputs.

\subsection{Tail energy and the limit of finite-subspace generation}
\begin{corollary}[Exact optimal truncation error]\label{cor:tail}
For any $\mu\in\Ptwo(H)$,
\begin{equation}\label{eq:tail-optimum}
 \inf_{\eta\in\Ptwo(H_m)}W_2^2(\mu,\eta)
 =W_2^2(\mu,(P_m)\push\mu)=\tau_m^2.
\end{equation}
If a centered Gaussian law has covariance eigenvalues $\lambda_k$ and $P_m$ retains its first $m$ eigenvectors, then $\tau_m^2=\sum_{k>m}\lambda_k$. If $\lambda_k\sim c k^{-2s}$ with $s>1/2$, then
\begin{equation}
 \tau_m\sim\sqrt{\frac{c}{2s-1}}\,m^{1/2-s}.
\end{equation}
\end{corollary}
\begin{proof}
Use the exact identity in \eqref{eq:pythagorean-certificate} and minimize its projected term. The spectral formula follows from Parseval's identity and the asymptotic from the integral test.
\end{proof}
The restriction in \eqref{eq:tail-optimum} is that outputs lie in the linear space $H_m$. It is not a general lower bound based only on latent dimension. If the full Gaussian has infinitely many strictly positive eigenvalues, its law and its finite projection are mutually singular even though their $W_2$ distance tends to zero. Accuracy in total variation, relative entropy, and Wasserstein distance must therefore be distinguished.
